% =============================================================================
%                   CHAPTER 2: LITERATURE REVIEW
% =============================================================================
\chapter{Literature Review}

\section{Background Study}

Video action recognition has evolved from hand-crafted spatiotemporal feature extractors to deep convolutional neural networks and self-supervised Vision Transformers. Early approaches relied on dense visual trajectory features such as Histogram of Oriented Gradients (HOG), Histogram of Optical Flow (HOF), and Improved Dense Trajectories (iDT), which captured local motion dynamics but struggled with long-range temporal modeling. The application of deep learning to video introduced 3D CNN architectures, including C3D and Inflated 3D ConvNets (I3D), which extended 2D kernels into 3D space to learn joint spatiotemporal features. I3D demonstrated that inflating ImageNet-pretrained weights enabled strong action classification performance on Kinetics and HMDB51. Further architectural refinements, such as SlowFast Networks, introduced dual-rate sampling pathways to decouple spatial semantics from rapid temporal motion.

While supervised video models achieve high accuracy, they depend heavily on expensive manual annotations. Self-supervised learning (SSL) resolves this limitation by learning representations from unlabelled video sequences via predictive pretext tasks. Generative video models, such as VideoMAE, train Vision Transformers (ViTs) to reconstruct masked pixel regions at high masking ratios (up to 90\%). Although VideoMAE achieves strong data efficiency, pixel-space reconstruction causes models to allocate representational capacity to high-frequency visual textures rather than high-level semantic action dynamics.

To overcome pixel-reconstruction overhead, Joint Embedding Predictive Architectures (JEPA) perform non-generative prediction directly in a learned latent feature space. \citet{assran2023ijepa} introduced Image-JEPA (I-JEPA), demonstrating that predicting missing visual target embeddings from context patches produces highly semantic representations. Extending JEPA to temporal sequences, \citet{bardes2024vjepa} proposed V-JEPA, which predicts spatiotemporal feature blocks in latent space. \citet{assran2025vjepa2} scaled video latent prediction with {\vjepaa{}}, a 654-million parameter ViT-Huge encoder trained on extensive video datasets without pixel reconstruction or language supervision. The V-JEPA~2 training objective performs \emph{mask-denoising prediction in representation space}: given two views $x$ and $y$ of a video, where patches have been randomly dropped from $y$, the encoder $E_\theta(\cdot)$ extracts representations from the unmasked view and a predictor $P_\phi(\cdot)$ predicts the masked representations by minimizing the $L_1$ distance to the stop-gradient target encoder output (Equation~\ref{eq:vjepa2_obj} in Chapter~3). The target encoder weights $\bar{\theta}$ are updated as an Exponential Moving Average of the context encoder (Equation~\ref{eq:vjepa_ema}), and a learnable mask token $\Delta_y$ indicates the positions of dropped patches. Both the encoder and predictor are parameterized as Vision Transformers with Rotary Position Embedding (RoPE).

Most recently, \citet{murlabadia2026vjepa21} introduced V-JEPA 2.1, which refines the training objective with a \emph{dense per-token prediction loss} (Equation~\ref{eq:vjepa21_obj}) and hierarchical intermediate supervision across encoder layers. Rather than computing a single global loss, V-JEPA~2.1 applies supervision individually to each masked spatiotemporal token, combined with a distance-weighted context loss on visible tokens (Equation~\ref{eq:vjepa21_context}). This dense formulation, combined with multi-layer intermediate supervision across encoder depth (Equation~\ref{eq:vjepa_multilayer}), significantly improves spatial and temporal representation grounding. The complete mathematical formulation of both V-JEPA~2 and V-JEPA~2.1 objectives is presented in Section~3.1.

Parallel to self-supervised feature learning, Vision-Language Models (VLMs) have emerged as powerful multimodal foundation models. As surveyed by \citet{bordes2024vlm} and \citet{zhang2023vlmsurvey}, VLMs connect visual encoders with autoregressive language decoders via visual instruction tuning. State-of-the-art vision-language models, such as \qwenvl{} \citep{bai2025qwen2vl} and Qwen3-VL \citep{qwen2025qwen3vl}, feature dynamic spatiotemporal tokenization, native multi-frame processing, and long-context multimodal understanding. To enable practical deployment under strict hardware resource boundaries, \citet{pathak2026efficientvlms} reviewed quantization, pruning, and low-bit inference techniques for large vision-language models.

The foundational computational mechanism underlying both self-supervised video backbones and vision-language architectures is the Transformer \citep{vaswani2017attention}. By replacing sequential recurrence with multi-head self-attention, the Transformer allows all positions in a sequence to interact directly via scaled dot-product attention:
\begin{equation}
\text{Attention}(\mathbf{Q}, \mathbf{K}, \mathbf{V}) = \text{softmax}\!\left(\frac{\mathbf{Q}\mathbf{K}^\top}{\sqrt{d_k}}\right)\mathbf{V}
\label{eq:self_attention}
\end{equation}
where queries $\mathbf{Q}$, keys $\mathbf{K}$, and values $\mathbf{V}$ are linear projections of the input sequence, and $\sqrt{d_k}$ stabilizes gradients. Each Transformer block pairs multi-head self-attention with a feed-forward network, stabilized by residual connections \citep{he2016resnet} and layer normalization \citep{ba2016layernorm}.

The Vision Transformer (ViT) \citep{dosovitskiy2021vit} adapted this architecture to visual data by dividing an image into non-overlapping spatial patches and treating them as visual tokens with added positional encodings. In modern video understanding, this paradigm is extended to 3D spatiotemporal tubelets, enabling models such as V-JEPA~2.1 and the visual front-end of \qwenvl{} to jointly learn spatial context and temporal motion dynamics without relying on rigid convolutional inductive biases.

\section{Literature Review}

The progression from Image-JEPA to V-JEPA~2.1 reveals a systematic trajectory of architectural refinement aimed at producing increasingly expressive latent representations. \citet{assran2023ijepa} first established the foundational principle that predicting target patch representations from context patches---without pixel-level reconstruction---produces semantically richer embeddings than generative or contrastive alternatives. Their Image-JEPA demonstrated that a non-generative latent predictive objective avoids wasting model capacity on perceptually irrelevant high-frequency detail, while simultaneously eliminating the dependency on hand-crafted data augmentations inherent to contrastive learning pipelines.

Building on this image-level foundation, \citet{bardes2024vjepa} extended the joint-embedding predictive paradigm to the spatiotemporal domain with V-JEPA, introducing masked spatiotemporal block prediction in video feature space. The critical insight was that temporal coherence across video frames provides an especially rich supervisory signal for latent prediction, as adjacent frames share semantic content while exhibiting continuous motion dynamics. \citet{assran2025vjepa2} then scaled the V-JEPA framework significantly with V-JEPA~2, employing a ViT-Huge backbone with 654 million parameters trained on large-scale unlabelled video data. This scaling demonstrated that increasing both model capacity and training data volume within the JEPA paradigm yields representations that transfer competitively across diverse downstream tasks.

The most recent architectural advancement, V-JEPA~2.1 \citep{murlabadia2026vjepa21}, introduced two critical innovations: a dense per-token prediction loss that supervises each masked spatiotemporal token individually, and multi-layer intermediate supervision applied across encoder depth. These modifications fundamentally change the nature of learned representations by ensuring that intermediate encoder layers also acquire strong semantic features, rather than concentrating representational quality only in the final layer. This property is especially relevant to the present study, as it motivates the extraction of multi-layer features for the proposed HAF-STD Net architecture. In parallel, \citet{huang2026vjepa} proposed Variational JEPA, which reformulates latent prediction as a probabilistic inference problem by introducing distribution uncertainty modelling, demonstrating that JEPA-style architectures can be extended beyond deterministic prediction to support stochastic control and predictive world modeling.

Integrating JEPA principles into multimodal and language domains represents an active frontier of research. \citet{chen2026vljepa} proposed VL-JEPA, extending joint-embedding predictive learning to multimodal vision-language alignment without relying on autoregressive token generation. By performing cross-modal latent prediction between visual and textual representations, VL-JEPA demonstrated that the non-generative JEPA objective can bridge modality gaps while preserving the computational advantages of avoiding pixel or token reconstruction. \citet{huang2025llmjepa} combined Large Language Models with JEPA, demonstrating latent-space predictive loss objectives for language understanding. Their LLM-JEPA framework showed that large language models benefit from predictive representation learning in ways that complement standard autoregressive training. \citet{liu2026jepareasoner} presented JEPA-Reasoner, decoupling latent-space reasoning from discrete token generation to enable unconstrained reasoning paths, while \citet{zhang2026thinkjepa} introduced ThinkJEPA, empowering latent world models with large vision-language reasoning capabilities. Collectively, these works indicate a growing trend toward unifying self-supervised prediction with language grounding, though none directly evaluate how these latent representations compare against instruction-tuned VLMs for discriminative video understanding tasks.

Beyond visual recognition, JEPA architectures serve as foundational world models for physical reasoning and decision-making. \citet{bagatella2025tdjepa} developed Temporal Difference JEPA, proving that latent-predictive representations enable zero-shot reinforcement learning across complex control environments. By combining temporal difference learning with JEPA's latent prediction, TD-JEPA demonstrated that forward predictive representations in a learned feature space encode sufficient dynamical information for effective policy transfer without environment-specific fine-tuning. \citet{terver2026physicaljepa} investigated the core drivers of success in physical planning using joint-embedding predictive world models, demonstrating strong generalization in physical simulation tasks. Their analysis identified that the latent prediction objective, rather than model scale alone, is the primary contributor to planning performance, providing evidence that JEPA representations capture causal-temporal structure that is qualitatively different from the associative representations learned through vision-language pretraining.

A fundamental distinction exists between the representation learning strategies of JEPA-based self-supervised models and instruction-tuned vision-language models. Self-supervised JEPA encoders such as V-JEPA~2.1 \citep{murlabadia2026vjepa21} learn visual representations through prediction in latent space, entirely without textual supervision. Their learned features are inherently visual-geometric, capturing motion trajectories, spatial configurations, and temporal dynamics as distributed patterns within a continuous embedding space. In contrast, vision-language models such as \qwenvl{} \citep{bai2025qwen2vl, qwen2025qwen3vl} acquire visual representations through cross-modal alignment with natural language, producing features that are linguistically grounded and semantically structured according to textual category boundaries. As discussed in the comprehensive survey by \citet{bordes2024vlm}, this linguistic grounding provides VLMs with strong zero-shot generalization capabilities, since visual concepts are inherently linked to compositional language descriptions. Recent efforts have sought to bridge the gap between linguistic and geometric representations; notably, \citet{lin2026qwen3d} introduced Qwen-3D, a generalist 3D vision-language model built upon the Qwen backbone that augments visual tokens with 3D Rotary Positional Embeddings and operates attention directly in 3D scene space, enabling scalable cross-view and temporal reasoning through explicit geometric awareness rather than purely linguistic alignment. However, \citet{zhang2023vlmsurvey} noted that vision-language alignment may bias representations toward semantically salient concepts that are easily described in natural language, potentially sacrificing fine-grained visual distinctions that lack clear linguistic labels.

The efficiency dimension further differentiates these paradigms. V-JEPA~2.1, with 654M parameters, achieves its representational power through architectural specialization and self-supervised pretraining on video data alone, whereas \qwenvl{} requires 7B parameters and massive multimodal corpora to achieve general-purpose visual understanding. \citet{pathak2026efficientvlms} surveyed quantization and pruning strategies for deploying large VLMs under hardware constraints, highlighting that the parameter efficiency of specialized self-supervised encoders represents a significant practical advantage for resource-constrained deployment scenarios.

Downstream representation evaluation depends heavily on task granularity. The FineGym benchmark establishes that gymnastics event and sub-action classification requires capturing subtle kinematic body movements across visually similar apparatus environments. HMDB51 contains 6,754 clips across 51 diverse action categories, with supervised baselines ranging from 57.2\% (iDT) to 80.7\% (I3D) and 73.3\% (VideoMAE).

The choice of downstream classifier significantly influences the evaluation of frozen representations. This study employs a diverse suite of classifiers spanning classical machine learning and modern neural network architectures. Support Vector Machines with Radial Basis Function kernels \citep{cortes1995svm} project features into high-dimensional spaces via the kernel trick, enabling non-linear decision boundaries and serving as strong baselines for feature evaluation. Random Forests \citep{breiman2001randomforests} construct ensembles of decision trees through bagging and random feature subsampling, providing robustness against overfitting and insensitivity to feature scaling. Logistic Regression \citep{cox1958logistic} applies a linear classifier with L2 regularization, serving as a linear separability probe that reveals whether learned representations are directly discriminative without non-linear transformation.

Among neural network classifiers, Multi-Layer Perceptrons (MLPs) \citep{rumelhart1986backprop} provide the foundational feedforward architecture with non-linear activation functions. Deeper variants with Batch Normalization \citep{ioffe2015batchnorm} and Dropout \citep{srivastava2014dropout} regularization enable training of deeper feedforward networks by stabilizing internal covariate shift and preventing co-adaptation of hidden units. Residual connections \citep{he2016resnet}, originally introduced for deep convolutional networks, enable gradient flow through skip connections of the form $\mathbf{h}_{l+1} = \mathbf{h}_l + \mathcal{F}(\mathbf{h}_l, \mathbf{W}_l)$, allowing substantially deeper architectures to be trained without degradation. One-dimensional convolutional networks \citep{lecun1998cnn} treat the feature dimension as a spatial axis, applying learned convolutional filters to capture local feature patterns and correlations. Finally, Transformer encoder blocks \citep{vaswani2017attention} employ multi-head self-attention to model global dependencies across feature dimensions, enabling each feature token to attend to all others through scaled dot-product attention.

\section{Research Gap and Motivation}

Despite significant advances in self-supervised video representation learning and multimodal foundation models, a comprehensive analysis of the existing literature reveals several critical research gaps that motivate the investigation undertaken in this thesis:

\textbf{Lack of Controlled Cross-Paradigm Comparison:} While self-supervised Joint Embedding Predictive Architectures (\vjepaa{}, \citealp{assran2025vjepa2}; V-JEPA 2.1, \citealp{murlabadia2026vjepa21}) and large vision-language models (\qwenvl{}, \citealp{bai2025qwen2vl, qwen2025qwen3vl}) have demonstrated remarkable individual capabilities, they have been evaluated exclusively within their own respective paradigms and benchmark suites. Self-supervised video encoders are typically assessed on masked-prediction benchmarks such as Kinetics-400 and Something-Something v2, whereas instruction-tuned vision-language models are primarily evaluated on visual question answering and text-generation tasks. To date, the literature lacks a direct, controlled empirical comparison of their frozen-feature representations under identical downstream classification conditions. Consequently, it remains an open question whether the visual-geometric features learned through latent-space predictive self-supervision or the linguistically grounded features acquired through multimodal instruction tuning provide superior transferability for video action understanding.

\textbf{Disparity Across Task Granularities:} Existing literature on video representation learning predominantly evaluates models on coarse-grained action recognition benchmarks such as HMDB51, where actions are characterized by distinct scenes, objects, and broad human movements. In contrast, fine-grained action recognition, exemplified by gymnastic routines in FineGym, requires distinguishing subtle kinematic differences, posture transitions, and limb trajectories in visually homogeneous environments. Whether the semantic abstractions of vision-language models or the temporal-predictive representations of JEPA encoders generalize more effectively across disparate levels of task granularity has not been systematically investigated.

\textbf{Under-Utilization of Multi-Layer Depth and Frequency Information:} Standard probing methodologies across both representation paradigms rely almost universally on simple global average pooling applied to the terminal encoder layer, followed by a linear classifier or shallow multi-layer perceptron. This conventional approach discards two vital sources of discriminative information. First, as established by \citet{murlabadia2026vjepa21}, intermediate layers in deep video encoders encode rich hierarchical features that capture complementary spatiotemporal abstractions across network depth. Second, global spatial and temporal pooling collapses temporal dynamics into static summary statistics, failing to capture high-frequency motion transients and multi-scale temporal variations. Existing downstream probing architectures provide no specialized mechanism to harness multi-layer depth representations alongside temporal frequency decomposition.

\textbf{Motivation and Contributions of this Study:} Motivated by these identified gaps, this thesis conducts the first rigorous, controlled empirical comparison between frozen representations extracted from V-JEPA 2.1 and \qwenvl{} across both fine-grained (FineGym) and coarse-grained (HMDB51) action recognition benchmarks. Both foundation models are systematically evaluated across a comprehensive suite of 11 diverse baseline classifiers, establishing fair and reproducible benchmark performance under identical protocol conditions. Furthermore, to address the architectural limitations of conventional probing, this study designs, implements, and evaluates the Hierarchical Adaptive Frequency Spatiotemporal Depth Network (HAF-STD Net), a novel classification architecture that leverages discrete wavelet multi-frequency decomposition, multi-depth feature extraction, and adaptive routing to maximize downstream spatiotemporal discrimination.
